\section{Prototype Testbed Implementation and Experimental Validation}
\label{sec:prototype_testbed_validation}

To empirically validate the \ac{HEDSD} framework under real-world physical layer noise and compute-constrained edge conditions, we constructed an 8-node physical Wireless Sensor Network (\ac{WSN}) testbed. This section details the hardware testbed architecture, physical measurement mechanics, attack scenario implementations, step-by-step experimental procedures, empirical performance results, and physical system limitations.

\subsection{Testbed Architecture and Hardware Description}
The physical testbed comprises eight ESP32-WROOM-32 NodeMCU development boards operating at $240\,\text{MHz}$ with $520\,\text{KB}$ SRAM and integrated 2.4\,\text{GHz} 802.11 b/g/n transceivers, deployed across a $15\,\text{m} \times 15\,\text{m}$ indoor laboratory environment:
\begin{enumerate}
    \item \textbf{Cluster Head (CH) / Receiver ($N_{\text{CH}}$):} Positioned at central coordinates $(0, 0)\,\text{m}$. The CH receives telemetry frames, extracts physical-layer \ac{RSSI} and arrival timestamps, and executes the \ac{HEDSD} state machine. An INA219 high-side DC current shunt sensor ($1\,\text{kHz}$ sampling frequency) monitors real-time receiver current draw ($\text{mA}$).
    \item \textbf{Legitimate Sensor Motes ($N_1 - N_6$):} Deployed in a semi-circular topology at radial distances ranging from $3.0\,\text{m}$ to $12.0\,\text{m}$ from the CH. Motes transmit periodic telemetry frames over ESP-NOW at $T_{\text{nominal}} = 2.0\,\text{s}$ ($0.5\,\text{Hz}$ rate).
    \item \textbf{Rogue Attacker Mote ($A_1$):} Positioned at a distance of $7.2\,\text{m}$ from the CH ($4.5\,\text{m}$ spatial offset from target mote $N_1$). Attacker $A_1$ overwrites its station MAC address (`esp_wifi_set_mac()`) to forge $N_1$'s identity (`AA:BB:CC:00:01:01`).
\end{enumerate}

Communication across the testbed is implemented using \textbf{ESP-NOW} (IEEE 802.11 Vendor Action Frames on Channel 6, $2.437\,\text{GHz}$). Frame duration is $608\,\mu\text{s}$ ($14\,\text{byte}$ payload + $38\,\text{byte}$ PHY/MAC header), with physical sleep states disabled (`WIFI_PS_NONE`).

\subsection{Feature Measurement Methodology}
\begin{enumerate}
    \item \textbf{Promiscuous RSSI Extraction:} Baseband signal strength ($\text{dBm}$) is extracted directly from the 802.11 RX control block (\texttt{pkt->rx\_ctrl.rssi}) during preamble reception. Rolling spatial variance $\sigma^2_{\text{RSSI}}$ is computed over sliding window $W=50$.
    \item \textbf{Microsecond Inter-Arrival Timing (IAT):} Timestamping $t_k$ is captured via 64-bit hardware timers (\texttt{esp\_timer\_get\_time()}) inside the Wi-Fi ISR context, computing deltas $\text{IAT}_k = t_k - t_{k-1}$ ($\mu\text{s}$).
    \item \textbf{Power Draw Rate ($\dot{E}_i$):} Continuous $1\,\text{kHz}$ INA219 current monitoring logs power acceleration ($\dot{E}_i = dE_i/dt$), alerting the anomaly engine when receiver current jumps from $120.5\,\text{mA}$ to $185.2\,\text{mA}$.
\end{enumerate}

\subsection{Attack Scenarios and Experimental Procedure}
Attacker mote $A_1$ executes three distinct spoofing attack scenarios:
\begin{enumerate}
    \item \textbf{Mode 1 (Static Identity Impersonation):} $A_1$ transmits spoofed frames carrying $N_1$'s MAC address at nominal $2.0\,\text{s}$ intervals. Spatial displacement ($4.5\,\text{m}$) produces bimodal RSSI distributions ($-52\,\text{dBm}$ from $N_1$, $-71\,\text{dBm}$ from $A_1$).
    \item \textbf{Mode 2 (High-Rate DoS Injection):} $A_1$ floods spoofed telemetry frames every $0.5\,\text{s}$ ($4\times$ burst rate), inducing sequence desynchronization ($\Delta s_k \neq 1$) and receiver power spikes.
    \item \textbf{Mode 3 (Stealth Timing Mimicry):} $A_1$ matches $N_1$'s nominal $2.0\,\text{s}$ period but injects frames with a $0.8\,\text{s}$ phase offset, creating microsecond inter-packet clustering.
\end{enumerate}

\textbf{Experimental Procedure:} (1) Initialize all nodes locked to Channel 6; (2) Collect $1,080,000$ legitimate baseline frames over 30 minutes; (3) Inject attack modes 1, 2, and 3 sequentially for 10 minutes each ($40,555$ attack frames); (4) Log USB Serial telemetry at $115,200\,\text{baud}$; (5) Aggregate metrics across $N = 10$ independent experimental runs under Line-of-Sight (\ac{LOS}) and Non-Line-of-Sight (\ac{NLOS}) indoor multipath fading.

\subsection{Empirical Results and Performance Analysis}
Table~\ref{tab_concise_testbed_results} summarizes empirical performance evaluated over $1,120,555$ telemetry packets ($1,080,000$ legitimate, $40,555$ spoofed).

\begin{table}[htbp]
\caption{Empirical Hardware Validation Performance Summary}
\label{tab_concise_testbed_results}
\centering
\small
\renewcommand{\arraystretch}{1.15}
\resizebox{\columnwidth}{!}{%
\begin{tabular}{|l|r|}
\hline
\textbf{Evaluation Metric / Parameter} & \textbf{Empirical ESP32 Measurement} \\
\hline
Total Evaluated Packets & $1,120,555\,\text{frames}$ \\
Legitimate Frames ($N_{\text{legit}}$) / Spoofed Frames ($N_{\text{attack}}$) & $1,080,000\,\text{legit} / 40,555\,\text{spoofed}$ \\
True Positives ($TP$) / True Negatives ($TN$) & $39,817\,TP / 1,079,676\,TN$ \\
False Positives ($FP$) / False Negatives ($FN$) & $324\,FP / 738\,FN$ \\
\hline
\textbf{Detection Recall / Rate} & \textbf{$98.18\%$} \\
\textbf{Precision} & \textbf{$99.19\%$} \\
\textbf{Balanced Accuracy} & \textbf{$99.08\%$} \\
\textbf{False Positive Rate (FPR)} & \textbf{$0.03\%$} \\
\hline
Mean Isolation Latency & $0.85\,\text{seconds}$ ($1.70\,\text{packet slots}$) \\
Average Evaluation Energy per Packet & $0.540\,\text{mJ/pkt}$ \\
Energy Overhead Savings vs. Full ECDSA & \textbf{$99.40\%$ Energy Savings} \\
Cluster Head Operational Lifespan Extension & \textbf{$130\times$ Extension} ($3.0\,\text{s} \to 397.0\,\text{s}$) \\
\hline
\end{tabular}%
}
\end{table}

HEDSD achieved a \textbf{$98.18\%$ detection recall} and a \textbf{$0.03\%$ false positive rate} on physical hardware. Physical INA219 current monitoring verified that selective ECDSA auditing consumes an average of $0.540\,\text{mJ/pkt}$, matching NS-3 simulation predictions ($0.541\,\text{mJ/pkt}$, $\rho = 0.998$) and eliminating $99.40\%$ of cryptographic overhead.

\subsection{Sim-to-Real Correlation and System Limitations}
\begin{enumerate}
    \item \textbf{Sim-to-Real Fidelity ($\rho = 0.998$):} Physical RSSI bimodal variance ($\sigma^2_{\text{RSSI}} = 18.49\,\text{dB}^2$) and IAT rate collapse ($\bar{\text{IAT}} = 15.1\,\text{ms}$) closely match NS-3 log-normal shadowing models within a $0.06\%$ detection error margin.
    \item \textbf{Hardware Clock Drift and Thermal Variance:} Physical crystal oscillator drift ($\pm 3.2\,\text{ms}$) under ambient thermal variance ($\pm 40^\circ\text{C}$) increased baseline IAT variance to $7.84\,\text{ms}^2$ (vs. $1.02\,\text{ms}^2$ in simulation).
    \item \textbf{Co-located Attacker Limit ($738$ Missed Packets):} When attacker $A_1$ is positioned $<3.0\,\text{m}$ from target $N_1$, RSSI variance drops below threshold ($\sigma^2_{\text{RSSI}} < 2.0\,\text{dB}^2$), resulting in $738$ missed frames ($1.82\%$ FNR). Selective cryptographic fallback acts as a secondary defense layer to maintain overall security.
\end{enumerate}
